\chapter{Степень формы, клиффордов след и происхождение относительного веса}
\label{ch:v7-clifford-form-degree-weight-origin-gate}

\section{Два естественных счёта степени}

После запрета простого Real-полуследа остаётся проверить, не различаются ли
рёберная и физическая кривизны своей степенью. Однако в текущем произведённом
носителе обе являются пространственно-временными нуль-формами:
\begin{equation}
 \deg_M\mathfrak m_E=\deg_M\mathfrak m_V=0.
 \label{eq:v7-clifford-common-spacetime-degree}
\end{equation}
Во внутреннем исчислении они также возникают на одном чётном уровне:
$\mathfrak m_E=[d_Z,d_Z^*]+\text{фон}$ и
$\mathfrak m_V=(A^*A-A_0^*A_0)\oplus(AA^*-A_0A_0^*)$ являются
квадратичными двухшаговыми моментами. Поэтому ни естественная внешняя, ни
внутренняя степень пока не создаёт относительного коэффициента.

\section{Явный клиффордов контроль}

Возьмём четырёхмерные евклидовы матрицы
\begin{equation}
 \{\gamma^\mu,\gamma^\nu\}=2\delta^{\mu\nu}I_4
 \label{eq:v7-clifford-relations}
\end{equation}
и клиффордово представление вещественной двухформы
\begin{equation}
 c(F)=\frac12\sum_{\mu,\nu}\gamma^\mu\gamma^\nu F_{\mu\nu}
 =\sum_{\mu<\nu}\gamma^\mu\gamma^\nu F_{\mu\nu}.
 \label{eq:v7-clifford-twoform-map}
\end{equation}
Прямой расчёт даёт изометрию
\begin{equation}
 \frac14\operatorname{tr}_{\rm spin}\bigl(c(F)^*c(F)\bigr)
 =\sum_{\mu<\nu}|F_{\mu\nu}|^2
 =\frac12\sum_{\mu,\nu}|F_{\mu\nu}|^2.
 \label{eq:v7-clifford-hodge-isometry}
\end{equation}
Для скаляра аналогично
\begin{equation}
 \frac14\operatorname{tr}_{\rm spin}|sI_4|^2=|s|^2.
 \label{eq:v7-clifford-scalar-isometry}
\end{equation}
Следовательно, видимая половина в последней части
\eqref{eq:v7-clifford-hodge-isometry} лишь удаляет двойной подсчёт пар
$(\mu,\nu)$ и $(\nu,\mu)$. Она не является весом между двумя секторами.

\section{Следствие для текущего действия}

Поскольку оба текущих момента тензорно умножаются на тождество спинорного
блока, нормированный клиффордов след сохраняет прежнее отношение:
\begin{equation}
 \boxed{\beta_{\rm Clifford}=1,\qquad
 (n_-,n_0,n_+)_{0}=(21,0,6).}
 \label{eq:v7-clifford-derived-weight-signature}
\end{equation}
Целевой вакуум остаётся положительным, но нулевой сектор запускает двадцать
один отрицательный режим вместо семи.

Это согласуется с суперсвязностным разложением кривизны на нуль-, одно- и
двухформенные компоненты: относительные веса становятся геометрическими
только после фиксации представления и следовой метрики; общий формализм сам
по себе допускает центральные веса
\cite{v7-clifford-roepstorff,v7-clifford-spectral-action,
v7-clifford-product-triples}.

\begin{theorem}[No-go степени формы и клиффордова следа текущего носителя]
В существующем произведённом носителе рёберная и физическая Hodge-кривизны
имеют одинаковую пространственно-временную степень и одинаковый чётный
внутренний статус. Нормированный четырёхмерный клиффордов след является
изометрией Hodge-нормы и сохраняет отношение $\beta=1$. Поэтому он не
выводит требуемое окно $0\le\beta<8/15$.
\end{theorem}

\begin{nogocriterion}[Запрет факториальной подгонки]
Нельзя объявлять только физический момент двухформой или смешивать
$\sum_{\mu<\nu}$ с $\sum_{\mu,\nu}$ после просмотра гессиана. Это меняет
представленное исчисление либо соглашение о компонентах только в одном
секторе и эквивалентно скрытому относительному весу.
\end{nogocriterion}

\begin{protocol}
\begin{enumerate}
 \item пространственно-временная степень $\mathfrak m_E$: \textbf{ноль};
 \item пространственно-временная степень $\mathfrak m_V$: \textbf{ноль};
 \item внутренний уровень обоих моментов: \textbf{чётный квадратичный};
 \item клиффордова изометрия двухформ: \textbf{проверена};
 \item фактор $1/2$ антисимметричных индексов: \textbf{не секторный вес};
 \item выведенное отношение: \textbf{$\beta=1$};
 \item нулевая сигнатура: \textbf{$(21,0,6)$};
 \item требуемое окно: \textbf{не достигнуто};
 \item итог: \textbf{текущий маршрут степени формы закрыт};
 \item следующий гейт: \textbf{кратности единого неприводимого носителя}.
\end{enumerate}
\end{protocol}

Машинный сертификат сохранён в
\path{s2t_v7_clifford_form_degree_weight_origin_gate_results.json}.

\begin{thebibliography}{9}
\bibitem{v7-clifford-roepstorff}
G.~Roepstorff,
\emph{Superconnections and the Higgs Field}, arXiv:hep-th/9801040.
\bibitem{v7-clifford-spectral-action}
A.~H.~Chamseddine, A.~Connes,
\emph{The Spectral Action Principle},
Commun. Math. Phys. 186 (1997), 731--750; arXiv:hep-th/9606001.
\bibitem{v7-clifford-product-triples}
S.~Brain, B.~Mesland, W.~D.~van Suijlekom,
\emph{Gauge Theory for Spectral Triples and the Unbounded Kasparov Product},
J. Noncommut. Geom. 10 (2016), 135--206; arXiv:1306.1951.
\end{thebibliography}